%-------------------------------------------------------------------------------
%	SECTION TITLE
%-------------------------------------------------------------------------------
\cvsection{Experience}


%-------------------------------------------------------------------------------
%	CONTENT
%-------------------------------------------------------------------------------
\begin{cventries}

%---------------------------------------------------------
  \cventry
    {Senior Scientist (Oberassistent)} % Job title
    {ETH Zurich - Department of Health Science and Technology, Institute of Robotics and Intelligent Systems, Sensory Motor Systems Lab} % Organization
    {Zurich, Switzerland} % Location
    {January 2018 - present} % Date(s)
    {
      \begin{cvitems} % Description(s) of tasks/responsibilities
        \item {Project lead in multiple engineering projects on optimization of Wheelchairs and Sport Climbing, including:}
        \item{Development of a physical wheelchair simulator that allows for the rapid evaluation of different seat positions in order to optimize ergonomics in everyday wheelchairs and aerodynamic performance in Paralympic racing wheelchairs. see \url{https://youtu.be/emKbxVdJjWc?t=85}}
        \item{Development of a measurement device for finger forces in sport climbers with visual real time feedback. Now used for regular, self guided assessment by the Swiss National and Olympic Climbing team.}
        \item{Actuation interface for a robotic bed that prevents sleep apnea.}
        \item{Key technologies used for these projects:\\C, C++, ARM, freeRTOS, ROS, Python 3.7, pygqt, Jupyter, MATLAB, Solidworks (CAD)}
        \item{Supervisor for ten master and bachelor thesis, six interns in engineering, and one senior engineer.}
      \end{cvitems}
    }
  
  %\cventry
  %  {Post Doc} % Job title
  %  {ETH Zurich - Department of Health Science and Technology, Institute of Robotics and Intelligent Systems, Sensory Motor Systems Lab} % Organization
  %  {Zurich, Switzerland} % Location
 %   {January 2018 - Dec 2019} % Date(s)
 %   {
 %     \begin{cvitems} % Description(s) of tasks/responsibilities
 %       \item {Optimization of the biomechanics in wheelchair propulsion: Numerically controlled, human in the loop, optimization of seating position to increase propulsion efficiency and prevent injuries, through a robotically controlled wheelchair ergometer with force sensing, kinematics tracking and motor controlled simulation of wheelchair dynamics}
  %      \item{Biomechanically modeling and analysis, training, assessment and rehabilitation in rock climbing: Automatized diagnostics, adaptive training and rehabilitation of the climbers fingers through force sensing combined with low latency visual feedback}
   %   \end{cvitems}
    %}
  
  \cventry
    {Research assistant} % Job title
    {ETH Zurich - Department of Health Science and Technology, Institute of Robotics and Intelligent Systems, Rehabilitation Engineering Lab} % Organization
    {} % Location
    {April 2012 - March 2017} % Date(s)
    {
      \begin{cvitems} % Description(s) of tasks/responsibilities
        \item {Research focus on upper-limb rehabilitation and assessment systems for patients with neurological injuries such as stroke and spinal cord injury \url{https://www.youtube.com/watch?v=ROiZXqdWeVk&ab_channel=RELABETHZ}}        
        \item {Development of ArmeoSenso, an affordable home therapy system with a touchscreen-based user interface in {Unity3D} and {C\#}}
        \item{Development of a novel, quaternion based, {sensor fusion} method for wearable, inertial measurement unit ({IMU}) based, low-latency motion capture of the patient’s arm, that is robust against drift and magnetic field disturbances, through classification of arm poses with suitable magnetic field}
        \item{Motion capture validated through {optical, marker-based tracking} (NaturalPoint Optitrack)}
        \item{Therapy consists of 3D movement exercises in the form of motivating and rewarding games as well as automatic diagnostics of the patient’s ability with resulting adaptation of the exercises}
        \item{Collaboration with clinicians and industry}
        \item {Clinical evaluation through a six week feasibility trial (NCT02098135) where patients trained with the system at home, on their own} 
        \item{Statistical analysis of the therapy efficacy and assessment validity}
        %\item{Supervisors: Prof. Roger Gassert, Dr. Olivier Lambercy, Prof. MD Andrease Luft}
        \item{Video summary: \url{www.relab.ethz.ch/research/past-research-projects/sensor-based_home_therapy_system.html}}
        \item{The system is now commercialized by Hocoma AG\ \url{www.hocoma.com/solutions/armeo-senso}}
        \item{Additional project supervision on classification of activities of daily living combined with encouraging feedback through wearable sensors, now Yband Therapy AG \url{ybandtherapy.com}}
      \end{cvitems}
    }

%---------------------------------------------------------
  \cventry
    {Master Thesis} % Job title
    {German Aerospace Center (DLR) - Institute for Robotics and Mechatronics} % Organization
    {Oberpfaffenhofen, Germany} % Location
    {Nov. 2010 - May 2011} % Date(s)
    {
      \begin{cvitems} % Description(s) of tasks/responsibilities
      	\item{Thesis: "Identification of Dynamic Parameters for a Biped Walking Machine Based on Foot- Ground Contact Force/Torque and Joint Torque Sensing"}
        \item {The humanoid robot consists of two DLR lightweight robots (commercially available as KUKA LBR iiwa) with force/torque sensors in the feet, controlled through {MATLAB \& SIMULINK}}
        \item{Dynamic, non-linear 18 DOF rigid body model derived with MAPLE {symbolic math}}
        \item {Feasible trajectory planing with {rapidly exploring randomized trees (RRTs)}} 
        \item{Selection of a set of trajectories, optimized for maximum parameter excitation and minimal experiment duration, through genetic algorithm}
        \item {Identification of inertial parameters through linear regression, validated through improved center of pressure prediction at the feet}
        \item{Supervisor: Dr. Christian Ott}
      \end{cvitems}
      }

%---------------------------------------------------------
  \cventry
    {Research Internship} % Job title
    { } % Organization
    { } % Location
    {Sep. 2009 - Oct. 2009} % Date(s)
    {
      \begin{cvitems} % Description(s) of tasks/responsibilities
        \item {Non-linear trajectory optimization and simulation for a robot arm with adaptive joint stiffness
}
        \item {Supervisor: Dr. Alin Albu-Schäffer}        
      \end{cvitems}
    }

%---------------------------------------------------------
  \cventry
    {Manufacturing Internship} % Job title
    { } % Organization
    { } % Location
    {Mar. 2005 - June 2005} % Date(s)
    {
      \begin{cvitems} % Description(s) of tasks/responsibilities        
        \item {CNC milling machines, lathes}
        \item{CAD and post-processing}
        \item {Digital position measurement based quality control}
      \end{cvitems}
    }

%---------------------------------------------------------
  \cventry
    {Visiting Student} % Job title
    {Massachusetts Institute of Technology (MIT) -  Active Adaptive Control Laboratory} % Organization
    {Boston, USA} % Location
    {Nov. 2009 - Apr. 2010} % Date(s)
    {
      \begin{cvitems} % Description(s) of tasks/responsibilities
        \item {Simulation and stability analysis of contact dynamics of quadrotors}
        \item {Supervisor: Dr. Anuradha Annaswamy}
      \end{cvitems}
      %\begin{cvsubentries}
      %  \cvsubentry{}{KNOX(Solution for Enterprise Mobile Security) Penetration Testing}{Sep. 2013}{}
      %  \cvsubentry{}{Smart TV Penetration Testing}{Mar. 2011 - Oct. 2011}{}
      %\end{cvsubentries}
    }

%---------------------------------------------------------
  \cventry
    {Student Thesis} % Job title
    {Technical University of Munich (TUM) -  Institute of  Automatic Control Engineering (LSR)} % Organization
    {Munich, Germany} % Location
    {Nov. 2008 - March. 2009} % Date(s)
    {
      \begin{cvitems} % Description(s) of tasks/responsibilities
        \item {Trajectory planning and control for the mobile base of a haptic teleoperation system to minimize end-effector vibrations}
        \item {Supervisor: Jens Hölldampf}
      \end{cvitems}
      %\begin{cvsubentries}
      %  \cvsubentry{}{KNOX(Solution for Enterprise Mobile Security) Penetration Testing}{Sep. 2013}{}
      %  \cvsubentry{}{Smart TV Penetration Testing}{Mar. 2011 - Oct. 2011}{}
      %\end{cvsubentries}
    }

%---------------------------------------------------------
  \cventry
    {Internship} % Job title
    {Technical University of Munich (TUM) -  Institute of Applied Mechanics (AMM)} % Organization
    { } % Location
    {Nov. 2006 - Oct. 2007} % Date(s)
    {
      \begin{cvitems} % Description(s) of tasks/responsibilities
        \item {Development on adaptive tessellation for a multi body dynamics visualization software (OpenMBV) in C++, Qt and OpenGL}
        \item {Computer aided design (3DS CATiA) of a joint calibrating mechanism for a biped robot}       
      \end{cvitems}
      %\begin{cvsubentries}
      %  \cvsubentry{}{KNOX(Solution for Enterprise Mobile Security) Penetration Testing}{Sep. 2013}{}
      %  \cvsubentry{}{Smart TV Penetration Testing}{Mar. 2011 - Oct. 2011}{}
      %\end{cvsubentries}
    }

%---------------------------------------------------------


\end{cventries}
